\documentclass[11pt]{article}
\usepackage[T1]{fontenc}
\usepackage[utf8]{inputenc}
\usepackage{lmodern,amsmath,amssymb,amsthm,booktabs,microtype}
\usepackage[letterpaper,margin=0.86in]{geometry}
\usepackage[colorlinks=true,urlcolor=blue,linkcolor=black]{hyperref}
\hypersetup{pdftitle={No universal cutoff for prescribed consecutive-pattern packings}}
\setlength{\parindent}{0pt}
\setlength{\parskip}{5pt plus 1pt minus 1pt}
\newcommand{\pos}[1]{(#1)_{+}}
\newcommand{\sourcepaper}{\href{https://arxiv.org/html/2610.29846v1\#S7}{Part I, \S7, item 8}}
\begin{document}
\begin{center}
{\Large\bfseries No universal cutoff for prescribed\\[2pt] consecutive-pattern packings}\par
\small Research note \enspace$\cdot$\enspace 18 November 2026
\end{center}

\textbf{Theorem.}
For every integer $h\ge1$, there are two permutation patterns of the same length and with the same prescribed overlap whose packing counts agree for the first $h$ terms and differ at term $h+1$. Both sequences are unbounded.

Blitvi\'c, Kammoun and Steingr\'imsson ask whether five occurrences suffice to distinguish any two different packing sequences. They report that, in their data through pattern length 12, pairs distinguishable within seven occurrences already differed within five. The theorem answers the universal-cutoff question in the negative. The construction also gives a lower bound for a cutoff allowed to depend on length. [\sourcepaper]

\subsection*{What a packing counts}
Let $p$ be a permutation of length $K$ and let $1\le m<K$. Write $N_{p,m}(j)$ for the number of permutations of length $K+(j-1)(K-m)$ in which each of the $j$ windows starting at $1+r(K-m)$, $0\le r<j$, has the same relative order as $p$. Other occurrences are allowed. Specifying a packing is equivalent to assigning ranks to all positions in a total order that respects the comparisons in every prescribed window. Each such order gives exactly one permutation, by assigning ranks from smallest to largest. [\href{https://arxiv.org/html/2610.29846v1\#S2.SS1}{Part I, Definition 2.6 and Remark 2.7}]

\subsection*{A family with an exact count}
Choose integers $d\ge1$ and $d+2\le w\le2d+1$. Write down blocks of three positions $(X_i,Y_i,Z_i)$ in order, starting at $i=0$. Assign them distinct integer keys
\[
    \kappa(X_i)=3i,\qquad \kappa(Y_i)=3(i+d)+1,\qquad \kappa(Z_i)=3(i+2d)+2 .
\]
The pattern $p_{w,d}$ is the list of ranks of the keys in the first $w$ blocks, read in position order. Its length is $K=3w$. Choose overlap $m=3(w-1)$, so that the window advances by one block. Put $(a)_+=\max(a,0)$. We claim that for every $j\ge1$,
\begin{equation}
   N_{p_{w,d},m}(j)=2^{\pos{w+j-2d-2}}.\label{eq:counts}
\end{equation}

\textit{Proof.}
There are $L=w+j-1$ blocks in the full permutation. Shifting the block index by one adds 3 to every key, so sorting all $3L$ positions by key gives one valid order. Two positions appear together in some prescribed window exactly when their block indices differ by at most $w-1$, since every run of at most $w$ consecutive blocks fits in a length-$w$ window. Thus every such pair has its key order forced.

In particular, each of $X_0,\ldots,X_{L-1}$, $Y_0,\ldots,Y_{L-1}$ and $Z_0,\ldots,Z_{L-1}$ is an increasing chain. Use the $Y$ chain to divide the order into gaps. If $i\le d$, then $X_i<Y_0$ is forced. For $i>d$, the forced comparisons are
\[
       Y_{i-d-1}<X_i<Y_{i-d}.
\]
If $i+d\ge L-1$, then $Y_{L-1}<Z_i$ is forced. For $i+d<L-1$ they are
\[
       Y_{i+d}<Z_i<Y_{i+d+1}.
\]
In each case the displayed comparisons follow from the keys and involve block distances at most $d+1\le w-1$. This pins every outer position to a particular gap, including the two end gaps.

The gap before $Y_0$ contains only $X$ positions, and the gap after $Y_{L-1}$ contains only $Z$ positions. Their orders are fixed. The only candidates in the gap $(Y_k,Y_{k+1})$ are
\[
                Z_{k-d} \quad\text{and}\quad X_{k+d+1},
\]
each present when its index lies in $[0,L-1]$. When both are present, their block indices differ by $2d+1\ge w$. No window contains the pair. They are consecutive in the full key order, so swapping them preserves every window comparison. Swaps in distinct gaps can be chosen independently. Conversely, the gap assignments and the two end chains fix every comparison except these choices.

Both candidates exist exactly when $d\le k\le L-d-2$. There are $(L-2d-1)_+$ such gaps. Each contributes exactly two orders. Substituting $L=w+j-1$ proves \eqref{eq:counts}.\hfill$\square$

\subsection*{Why every proposed cutoff fails}
Given $h\ge1$, choose $w=h+4$ and use the offsets $d=h+1$ and $d'=h+2$. Both meet the inequalities of the construction. Formula \eqref{eq:counts} gives the two counts, for every $j$,
\[
      2^{\pos{j-h}} \qquad\text{and}\qquad 2^{\pos{j-h-2}}.
\]
They are both 1 for $1\le j\le h$. At $j=h+1$ the counts are 2 and 1. This proves the theorem. For length $K=3(h+4)$ these two patterns have common overlap $K-3$, so a distinguishing cutoff for that length must be at least $h+1=K/3-3$. The optimal length-dependent cutoff is a separate question.

\textbf{Example.} For the proposed cutoff of five, take $w=9$. Both patterns have length 27 and overlap 24. In the first pattern at $j=6$, the first swappable pair is $Z_0,X_{13}$. Their keys are 38 and 39. They lie thirteen blocks apart, and each window sees only nine blocks.
\begin{center}\small
\begin{tabular}{lrrrrrrrrr}
\toprule
Number of occurrences $j$ & 1&2&3&4&5&6&7&8&9\\
\midrule
Offset $d=6$ &1&1&1&1&1&\textbf{2}&4&8&16\\
Offset $d=7$ &1&1&1&1&1&\textbf{1}&1&2&4\\
\bottomrule
\end{tabular}
\end{center}

\textbf{Reproducible check.}
Run \texttt{python3 verify.py} with the accompanying file. It forms precedence constraints directly from the ranks in each window and counts total orders by recursively choosing the next available position. It reproduced the table and matched \eqref{eq:counts} in 754 cases covering $1\le d\le12$, all permitted $w$, and values of $j$ on both sides of the first split. The proof above applies to all parameter values.

\textbf{Status of the question.}
The authors pose the cutoff in question form. As of 18 November 2026, searches using the paper's authors, its arXiv identifier, and terms including \emph{cutoff}, \emph{five occurrences}, and \emph{initial agreement} turned up no earlier resolution of the universal assertion. The companion paper studies moment sequences, Nagy's construction answers item 6, and a separate public proof of the packing-maximum result explicitly distinguishes item 8 as outstanding. This is a bounded literature check. The proof here has not undergone peer review. [\sourcepaper, \href{https://arxiv.org/abs/2611.12829}{Part II}, \href{https://arxiv.org/abs/2611.06937}{Nagy}, \href{https://the-omega-institute.github.io/trureturing-mdbook/Blueprint/D5/S3/Combinatorics/Permutation/PackingMaximumRefutation.html}{packing-maximum proof}]
\end{document}
